% python_api.tex — direct object-level Python API usage (GitTree / GitRepo)

\chapter{Direct Python Object API}
\label{chap:python-api}

This chapter shows how to manipulate \cgspkg{} objects directly (without the
CLI wrapper), including loading a \texttt{.gts} snapshot, reconstructing the
runtime registry, mirroring identity into \texttt{GitTree}/\texttt{GitRepo},
and propagating a tag target.

For normal Multi-repo Sync usage, prefer the CLI lifecycle documented in the
README and Getting Started guide. The reference test topology is CGSil1
(\texttt{https://gitlab.com/CGS\_test/CGSil1}):
\texttt{initialise -> pull -> checkout/add/commit/push/tag -> freeze}.

The lower-level object/API lifecycle remains:
\texttt{load(.cgs) -> .gts LOADED -> expand(.cgs|.gts) -> .gts PENDING ->
validate(.cgs|.gts) -> .gts READY -> pull(.cgs|.gts) -> checkout(.gts) ->
add -> git(tree,"commit",msg) -> git(tree,"push") -> git(tree,"tag",name) ->
freeze}.

\section{From Empty Registry to READY via \texttt{.gts}}

\begin{lstlisting}[language=Python]
from pathlib import Path

from ComplexGitSync import GitRepo, GitTree, RefKind, TreeLifecycleState
from ComplexGitSync.orchestre import GtsDocument, build_registry_from_gts_document
from ComplexGitSync.git_tree import WorkingGitTree

# 1) Empty registry lifecycle
empty_registry = WorkingGitTree()
assert empty_registry.recompute_tree_state() == TreeLifecycleState.UNLOADED

# 2) Load the CGSil1 runtime snapshot (.gts) and rebuild the dependency registry
cgshome = Path("../..")
gts_doc = GtsDocument.from_toml(cgshome / ".cgitsync/state/CGSil1.gts")
registry = build_registry_from_gts_document(gts_doc)
assert registry.lifecycle_state == TreeLifecycleState.READY

# 3) Rebuild GitTree/GitRepo identity objects from discovered registry entries
tree = GitTree()
for entry in registry.values():
    tree.add_repo(
        GitRepo(
            project_owner_name=entry.project_owner_name or "owner",
            project_name=entry.project_name or entry.name,
            gitprovider=entry.gitprovider,
            access_protocol=entry.access_protocol,
            commit_sha=entry.commit_sha,
        )
    )
\end{lstlisting}

\section{Object-Level Tag Propagation}

\begin{lstlisting}[language=Python]
# Set the same tag target across all registry entries (parent -> leaves)
tree.propagate_tag(registry, "v1.2.3")
for entry in registry.values():
    assert entry.target_ref_kind == RefKind.TAG
    assert entry.target_ref_name == "v1.2.3"
\end{lstlisting}

This object-level flow is the same state model used by the CLI and
\texttt{ComplexGitSyncClient}, but exposed directly for advanced automation and
tests.

\section{READY Methods in the Client API}

The same lifecycle model is exposed by \texttt{ComplexGitSyncClient}.  In
practice, the workflow is:

\begin{itemize}
  \item \texttt{load(...)} is the step-1 name for direct source loading.
  \item \texttt{expand(...)} is the canonical step-2 name; it runs nested
    \texttt{.cgs} discovery from parents to leaves (recursive).
  \item \texttt{validate(...)} is the canonical step-3 name; it accepts both
    \texttt{.cgs} and \texttt{.gts} sources.
  \item \texttt{initialise(...)} bootstraps a tree and must finish in
    \texttt{READY}; \texttt{clone(...)} is the direct clone entry point for a
    \texttt{.cgs} source; \texttt{pull(...)} resynchronises an existing tree.
  \item \texttt{git(tree, command, *args)} is the unified step-5 interface for
    \texttt{"commit"}, \texttt{"push"}, and \texttt{"tag"} operations.  Each
    follows leaf-first ordering.
  \item \texttt{freeze(...)} emits the next persisted \texttt{.gts} state.
\end{itemize}

\begin{lstlisting}[language=Python]
from pathlib import Path

from ComplexGitSync import ComplexGitSyncClient

client = ComplexGitSyncClient()
cgspath = Path("../..")
cgshome = cgspath / "CGSil1"

# 1. load the CGSil1 test-case .cgs -> .gts LOADED
client.load("../CGSil1.cgs")

# 2. expand(.cgs/.gts) -> .gts PENDING
print(client.expand("../CGSil1.cgs"))

# 3. validate(.cgs/.gts) -> .gts READY
client.validate("../CGSil1.cgs")

# 4. clone(.cgs/.gts)
# Same default as the CLI: output_path defaults to CGSPATH=../..
client.initialise("../CGSil1.cgs")

# Equivalent explicit form:
client.initialise("../CGSil1.cgs", output_path=cgspath)

# Recovery path for partial/stale clone state:
client.clean_initialise_cgs("../CGSil1.cgs", output_path=cgspath)

# Cleanup only:
client.purge_cgs("../CGSil1.cgs", output_path=cgspath)

# 5. READY methods — unified git interface
registry = client.get_dependency_registry()
client.checkout("feature/my-branch")
client.add()
client.git(registry, "commit", "feat: update CGSil1 CGS#1")
client.git(registry, "push")
client.git(registry, "tag", "v1.2.3")

# 6. freeze -> next .gts id
client.freeze("release-2026.05")
\end{lstlisting}
